\subsection{Comparison with Prior Studies}
\label{sec:comparison_with_prior_studies}
\label{sec:comparison_prior}

Perbandingan empiris kerangka kerja usulan terhadap model terdahulu dan baseline internal pada dataset DemogPairs disajikan pada Table~\ref{tab:15}. Evaluasi distrukturkan ke dalam dua kelompok protokol partisi data, yaitu kelompok protokol Stratified dan kelompok protokol Subject-Disjoint. Studi seminal oleh Hupont dan Fern{\'a}ndez (2019)~\cite{ref25} menguji verifikasi pasangan identitas berbasis kurva Receiver Operating Characteristic (ROC) dan Area Under the Curve (AUC), serta tidak mengevaluasi klasifikasi 6-kelas interseksional. Pada kelompok protokol Stratified yang merepresentasikan komparasi kontekstual, Proposed Tri-Domain + SVM (Historical Baseline) mencatat performa tertinggi dengan Akurasi 93.70\%, Presisi 93.72\%, Recall 93.70\%, dan F1-Score 93.69\%. Capaian tersebut melampaui model terdahulu Putri et al. (2025)~\cite{ref19} (Akurasi 91.15\%) dan Putri et al. (2025)~\cite{ref20} (90.85\%). Selain itu, hasil ini juga mengungguli baseline komparasi ViT-Base~\cite{ref17} (85.40\%) dan ResNet-50~\cite{ref13} (81.25\%). Pada kelompok protokol Subject-Disjoint yang mengevaluasi generalisasi tanpa kebocoran identitas, Proposed Tri-Domain + LR memimpin tolok ukur dengan Akurasi 88.66\%, Presisi 88.75\%, Recall 88.66\%, dan F1-Score 88.68\%, disusul oleh Proposed Tri-Domain + SVM (88.61\%). Fusi tri-domain secara konsisten mengungguli konfigurasi single-domain dan dual-domain.

\begin{table*}[htbp]
\caption{Contextual and Controlled Comparison of the Proposed Framework on the DemogPairs Dataset.}
\label{tab:15}
\label{tab:XV}
\centering
\begin{tabular}{llccccc}
\toprule
Model & Feature Domain & Data Split Protocol & Accuracy (\%) & Precision (\%) & Recall (\%) & F1-Score (\%) \\
\midrule
Hupont \& Fern{\'a}ndez (2019)~\cite{ref25}$^*$ & Single (Face Verification) & Pairwise ROC Verification & -- & -- & -- & -- \\
\midrule
Baseline ResNet-50~\cite{ref13} & Single (Face) & Stratified & 81.25 & 81.30 & 81.25 & 81.20 \\
Baseline ViT-Base~\cite{ref17} & Single (Face) & Stratified & 85.40 & 85.45 & 85.40 & 85.38 \\
Putri et al. (2025)~\cite{ref19} & Dual (Face + Emotion) & Stratified & 91.15 & 91.20 & 91.15 & 91.12 \\
Putri et al. (2025)~\cite{ref20} & Dual (Face + Age) & Stratified & 90.85 & 90.90 & 90.85 & 90.82 \\
\textbf{Proposed Tri-Domain + SVM (Historical Baseline)} & \textbf{Tri (Face + Emotion + Age)} & \textbf{Stratified} & \textbf{93.70} & \textbf{93.72} & \textbf{93.70} & \textbf{93.69} \\
\midrule
Proposed Single-Domain + LR & Single (Face) & Subject-Disjoint & 85.28 & 85.35 & 85.28 & 85.27 \\
Proposed Single-Domain + SVM & Single (Face) & Subject-Disjoint & 85.74 & 85.80 & 85.74 & 85.73 \\
Proposed Dual-Domain + LR & Dual (Face + Emotion) & Subject-Disjoint & 87.27 & 87.35 & 87.27 & 87.28 \\
Proposed Dual-Domain + SVM & Dual (Face + Emotion) & Subject-Disjoint & 88.47 & 88.52 & 88.47 & 88.48 \\
\textbf{Proposed Tri-Domain + LR} & \textbf{Tri (Face + Emotion + Age)} & \textbf{Subject-Disjoint} & \textbf{88.66} & \textbf{88.75} & \textbf{88.66} & \textbf{88.68} \\
Proposed Tri-Domain + SVM & Tri (Face + Emotion + Age) & Subject-Disjoint & 88.61 & 88.66 & 88.61 & 88.63 \\
\bottomrule
\end{tabular}
\par\vspace{2pt}\raggedright\footnotesize $^*$(Catatan: Hupont \& Fern{\'a}ndez (2019) mengevaluasi tugas verifikasi pasangan identitas wajah berbasis kurva ROC/AUC, bukan tugas klasifikasi 6-kelas interseksional).
\end{table*}

Penetapan posisi riset ini menegaskan kualifikasi metodologis yang krusial antara kedua protokol evaluasi tersebut. Korpus DemogPairs orisinal yang diperkenalkan oleh Hupont dan Fern{\'a}ndez~\cite{ref25} sejatinya dirancang sebagai tolok ukur verifikasi identitas wajah berpasangan dan bukan klasifikasi atribut multi-kelas tertutup. Oleh sebab itu, penyelidikan ini bersama studi literatur terkait memelopori penyediaan baseline komprehensif untuk klasifikasi 6-kelas interseksional pada dataset tersebut. Hasil pada protokol Stratified diklasifikasikan murni sebagai komparasi kontekstual karena pemisahan acak tingkat citra membiarkan identitas subjek tumpang-tindih antara set latih dan uji, ditambah potensi variasi teknik prapemrosesan antarstudi. Perbedaan performa antara baseline historis Stratified (93.70\%) dan benchmark primer Subject-Disjoint (88.66\%) mengindikasikan adanya optimisme semu akibat kebocoran data identitas. Penerapan protokol ketat Subject-Disjoint 480 subjek latih dan 120 subjek uji held-out menjamin eliminasi kebocoran identitas secara mutlak, sehingga menghasilkan evaluasi performa yang andal dan mencerminkan kapasitas generalisasi sistem biometrik pada dunia nyata.
